\documentclass[11pt,a4paper]{article}

% ── Packages ─────────────────────────────────────────────────────────────────
\usepackage{amsmath,amssymb,amsthm,bm}
\usepackage[margin=2.5cm]{geometry}
\usepackage{hyperref}
\usepackage{parskip}
\usepackage{booktabs}

\newtheorem{theorem}{Theorem}

% ── Title ─────────────────────────────────────────────────────────────────────
\title{\textbf{Contraction Theory --- Derivation Notes}\\
       \large (Note 20 Feb 2026)}
\author{Dipankar Bhattacharya}
\date{March 2026}

\begin{document}
\maketitle
\tableofcontents
\newpage

% ─────────────────────────────────────────────────────────────────────────────
\section{Setup}

Consider a smooth nonlinear system
\begin{equation}
  \dot{\bm{x}} = \bm{f}(\bm{x}, t), \qquad \bm{x}\in\mathbb{R}^n.
  \label{eq:sys}
\end{equation}
Let $\bm{x}_1(t)$ and $\bm{x}_2(t)$ be \emph{two solutions} of \eqref{eq:sys}
starting from different initial conditions (ICs).
Define the \textbf{virtual displacement}
\[
  \delta\bm{x} = \bm{x}_1 - \bm{x}_2 \in \mathbb{R}^n,
\]
which satisfies the \textbf{variational dynamics} (linearised along any
trajectory):
\begin{equation}
  \delta\dot{\bm{x}} = \bm{F}(\bm{x},t)\,\delta\bm{x},
  \qquad
  \bm{F} = \frac{\partial \bm{f}}{\partial \bm{x}}.
  \label{eq:var}
\end{equation}

% ─────────────────────────────────────────────────────────────────────────────
\section{Metric and Transformed Displacement}

Introduce a uniformly positive-definite \textbf{metric} (coordinate
transformation)
\[
  \bm{\Theta}(\bm{x},t) \in \mathbb{R}^{n\times n}, 
  \qquad \bm{\Theta}^\top\bm{\Theta} = \bm{M}(\bm{x},t) \succ 0,
\]
and define the \textbf{transformed virtual displacement}
\begin{equation}
  \delta\bm{y} = \bm{\Theta}\,\delta\bm{x}.
  \label{eq:dy}
\end{equation}

The \textbf{squared length} (Lyapunov-like measure) is
\begin{equation}
  V = \delta\bm{y}^\top\delta\bm{y}
    = \delta\bm{x}^\top \bm{M}\,\delta\bm{x} \geq 0.
  \label{eq:V}
\end{equation}

% ─────────────────────────────────────────────────────────────────────────────
\section{Time Derivative of Squared Length}

Differentiating \eqref{eq:V} along \eqref{eq:var}:
\begin{align}
  \dot V
  &= \delta\dot{\bm{y}}^\top\delta\bm{y}
   + \delta\bm{y}^\top\delta\dot{\bm{y}} \notag\\
  &= 2\,\delta\bm{y}^\top\delta\dot{\bm{y}}. \label{eq:Vdot1}
\end{align}

From \eqref{eq:dy}: $\delta\dot{\bm{y}} = \dot{\bm{\Theta}}\,\delta\bm{x}
+ \bm{\Theta}\,\delta\dot{\bm{x}}
= \bigl(\dot{\bm{\Theta}} + \bm{\Theta}\bm{F}\bigr)\delta\bm{x}
= \bigl(\dot{\bm{\Theta}} + \bm{\Theta}\bm{F}\bigr)\bm{\Theta}^{-1}\delta\bm{y}.$

Define the \textbf{generalised Jacobian}
\begin{equation}
  \bm{K}(\bm{x},t)
  = \bigl(\dot{\bm{\Theta}} + \bm{\Theta}\bm{F}\bigr)\bm{\Theta}^{-1},
  \label{eq:K}
\end{equation}
so that $\delta\dot{\bm{y}} = \bm{K}\,\delta\bm{y}$, and \eqref{eq:Vdot1}
becomes
\begin{equation}
  \boxed{
  \dot V = 2\,\delta\bm{y}^\top \bm{K}\,\delta\bm{y}
         = \delta\bm{y}^\top\bigl(\bm{K} + \bm{K}^\top\bigr)\delta\bm{y}.
  }
  \label{eq:Vdot}
\end{equation}

% ─────────────────────────────────────────────────────────────────────────────
\section{Contraction Condition}

\begin{theorem}[Contraction]
The system \eqref{eq:sys} is \textbf{contracting} with respect to the metric
$\bm{\Theta}$ if there exists a constant $\lambda > 0$ such that
\begin{equation}
  \bm{K}_s \triangleq \frac{1}{2}\bigl(\bm{K} + \bm{K}^\top\bigr)
  \preceq -\lambda\,\bm{I}, \quad \forall\,\bm{x},\,t.
  \label{eq:cond}
\end{equation}
That is, the \emph{symmetric part} of $\bm{K}$ is uniformly negative definite.
\end{theorem}

\begin{proof}
From \eqref{eq:Vdot} and \eqref{eq:cond}:
\[
  \dot V \leq -2\lambda\,V.
\]
By Gr\"{o}nwall's lemma:
\[
  V(t) \leq V(0)\,e^{-2\lambda t},
\]
i.e.\
\begin{equation}
  \|\delta\bm{y}(t)\| \leq \|\delta\bm{y}(0)\|\,e^{-\lambda t}.
\end{equation}
Since $\bm{\Theta}$ is uniformly bounded and invertible, this gives
\textbf{exponential convergence} of all trajectories to each other:
\[
  \|\bm{x}_1(t) - \bm{x}_2(t)\| \to 0 \quad \text{exponentially}.
\]
\end{proof}

% ─────────────────────────────────────────────────────────────────────────────
\section{Special Case: Identity Metric}

With $\bm{\Theta} = \bm{I}$ the generalised Jacobian reduces to the ordinary
Jacobian $\bm{K} = \bm{F}$, and the contraction condition becomes:
\begin{equation}
  \bm{F}_s = \frac{1}{2}\bigl(\bm{F} + \bm{F}^\top\bigr) \preceq -\lambda\bm{I}.
  \label{eq:euclidean}
\end{equation}
This is the \emph{Euclidean} contraction condition.

% ─────────────────────────────────────────────────────────────────────────────
\section{Examples}

\subsection*{Example (a) --- Linear time-invariant system}

Let $\dot{\bm{x}} = A\bm{x}$, $A\in\mathbb{R}^{n\times n}$ constant.
With $\bm{\Theta} = \bm{S}$ (constant), $\bm{K} = \bm{S}A\bm{S}^{-1}$.\\
Choose $\bm{S}^\top\bm{S} = \bm{M}$ such that $\bm{K}+\bm{K}^\top \prec 0$:
\[
  \bm{S}A\bm{S}^{-1} + \bm{S}^{-\top}A^\top\bm{S}^\top \prec 0
  \;\Longleftrightarrow\;
  A^\top\bm{M} + \bm{M}A \prec 0 \quad (M = S^\top S).
\]
This is the standard Lyapunov stability condition — contraction subsumes it.

\subsection*{Example (b) --- Scalar nonlinear system}

Let $\dot x = f(x) = -x - x^3$.  Then $F = \partial f/\partial x = -1 - 3x^2$.
For all $x$, with $\bm{\Theta}=1$:
\[
  K_s = F = -1 - 3x^2 \leq -1 < 0.
\]
So $\lambda = 1$; the system is globally contracting and all trajectories
converge exponentially.

% ─────────────────────────────────────────────────────────────────────────────
\section{Summary}

\begin{center}
\small
\begin{tabular}{@{}ll@{}}
\toprule
\textbf{Quantity} & \textbf{Definition / Role} \\
\midrule
$\delta\bm{x}$ & Virtual displacement between two trajectories \\
$\bm{\Theta}$ & Uniformly invertible coordinate transformation (metric) \\
$\delta\bm{y} = \bm{\Theta}\,\delta\bm{x}$ & Transformed displacement \\
$V = \|\delta\bm{y}\|^2$ & Squared length (Lyapunov measure) \\
$\bm{K} = (\dot{\bm{\Theta}}+\bm{\Theta}\bm{F})\bm{\Theta}^{-1}$ & Generalised Jacobian \\
$\bm{K}_s \preceq -\lambda\bm{I}$ & Contraction condition ($\lambda>0$) \\
$\|\delta\bm{x}\|\to0$ exponentially & Conclusion (all-to-all convergence) \\
\bottomrule
\end{tabular}
\end{center}

\vspace{6pt}
\noindent
\textbf{Key insight:} A system is contracting if a metric $\bm{\Theta}$ can be
found such that the symmetric part of the generalised Jacobian $\bm{K}$ is
uniformly negative definite.  Unlike Lyapunov methods, this gives
\emph{any-to-any} trajectory convergence, not just stability of a single
equilibrium --- making it directly applicable to the cup-and-ball transport task
studied in Bazzi et al.\ (2018).

\end{document}
